% Laget automatisk av elevark/lag_elevark.py fra sim/e-over-m.html. Endringer her blir overskrevet.
\documentclass[11pt]{article}
\usepackage{elevark}
\begin{document}
\elevarkhode{Fysikk 2 · Magnetiske felt og krefter}{$e/m$ for elektronet}{https://lektorbo.no/sim/e-over-m.html}{Simuleringen}

I et fintrådrør blir elektroner akselerert gjennom spenningen $U$ og sendt inn i magnetfeltet fra to Helmholtz-spoler. Litt gass i røret lyser der elektronene går, så vi ser banen som en sirkel. Fra diameteren kan vi regne ut forholdet mellom ladningen og massen til elektronet.

\section*{Oppgaver}

Bruk $e = 1{,}602 \cdot 10^{-19}\,\mathrm{C}$ og $m_{\mathrm{e}} = 9{,}109 \cdot 10^{-31}\,\mathrm{kg}$ når du regner med kjente verdier.

\begin{oppgaveliste}
\oppgave \fase{gjett}{Gjett}Hva skjer med sirkelen hvis du øker spenningen $U$? Hva skjer hvis du øker strømmen $I$ i spolene? Test begge.
\svar
\oppgave \fase{test}{Regn}Hvor stor fart får elektronene med $U = 200\,\mathrm{V}$? Hvor mange prosent av lysfarten er det?
\svar
\oppgave \fase{test}{Regn}Regn ut magnetfeltet og diameteren på banen når $U = 200\,\mathrm{V}$ og $I = 1{,}50\,\mathrm{A}$. Sjekk med målestaven, og kryss av for å se diameteren som tall.
\svar
\oppgave \fase{forklar}{Forklar}Vis at $e/m = 2U/(B^{2}r^{2})$. Hvilke to formler må du kombinere?
\svar
\oppgave \fase{test}{Mål}Skjul tallet for diameteren. Velg $U = 250\,\mathrm{V}$ og $I = 1{,}40\,\mathrm{A}$, og les av diameteren på målestaven. Regn ut $e/m$. Hvor stort er avviket fra den kjente verdien?
\svar
\oppgave \fase{test}{Undersøk}Hold $U = 250\,\mathrm{V}$ fast, og mål diameteren for $I = 1{,}0$, 1,2, …, 2,0~A. Plott $1/r$ mot $I$. Hvorfor blir det en rett linje, og hvordan kan du finne $e/m$ fra stigningstallet?
\svar
\oppgave \fase{forklar}{Forklar}Den magnetiske kraften gjør at elektronene går i sirkel, men farten endrer seg ikke. Hvorfor? Hva ville skjedd med banen hvis elektronene kom inn litt på skrå i forhold til feltet?
\svar
\oppgave \fase{utfordring}{Utfordring}Et proton har like stor ladning som et elektron, men 1836 ganger så stor masse. Hvor stor ville radien blitt for et proton med samme $U$ og $B$? Kunne du brukt dette røret til å måle $e/m$ for protoner?
\svar
\end{oppgaveliste}

\end{document}
